\section{\texorpdfstring{A uniform heat estimate on \(\USp(2g)\)}%
  {A uniform heat estimate on USp(2g)}}
\label{app:heat-trace}

We prove the compact-group estimate used in \cref{sec:transfer}, including
the metric, Casimir, and Haar normalizations.  Throughout this appendix,
\(G_g=\USp(2g)\), \(m_g^{\mathrm{Haar}}(G_g)=1\), and the metric is
\eqref{eq:group-metric}.  If \(X_1,\ldots,X_{D_g}\) is an orthonormal basis
of left-invariant vector fields, where \(D_g=2g^2+g\), set
\begin{equation}
 \mathcal A_g=-\sum_{j=1}^{D_g}X_j^2,
 \qquad P_s=e^{-s\mathcal A_g}.
\label{eq:heat-generator}
\end{equation}
Thus the Euclidean model of this generator has coordinate variance \(2s\).
We write \(K_s(x)\) for the heat-kernel density relative to Haar
\emph{probability} measure, not relative to unnormalized Riemannian volume;
see also \cite{Hall1999}.

\subsection{Root coordinates and the heat trace}

On the standard maximal torus, write
\[
 X(x)=\operatorname{diag}(ix_1,\ldots,ix_g,-ix_1,\ldots,-ix_g).
\]
Direct multiplication gives
\begin{equation}
 -\frac12\operatorname{Tr}(X(x)X(y))=\sum_{i=1}^gx_iy_i.
\label{eq:torus-metric}
\end{equation}
The type-\(C_g\) positive roots are therefore expressed in the Euclidean
basis \(e_i\), and their half-sum is
\(\varrho=(g,g-1,\ldots,1)\).  Irreducibles are indexed by partitions
\(\gamma=(\gamma_1\ge\cdots\ge\gamma_g\ge0)\).  With the normalization
\eqref{eq:heat-generator}, the Casimir eigenvalue is exactly
\begin{equation}
 \kappa_\gamma
 =\|\gamma+\varrho\|^2-\|\varrho\|^2
 =\sum_{i=1}^g\gamma_i
       \bigl(\gamma_i+2(g-i+1)\bigr).
\label{eq:casimir}
\end{equation}
For example, at \(g=1\) the defining representation has
\(\kappa_{(1)}=3\), agreeing with
\(-\sum_{j=1}^3(i\sigma_j)^2=3I\) for the orthonormal basis
\(i\sigma_1,i\sigma_2,i\sigma_3\).  This check rules out a hidden factor of
two in the heat time.

Let \(d_\gamma\) be the representation dimension.  The type-\(C_g\) Weyl
formula in these coordinates is
\begin{equation}
\begin{split}
 d_\gamma={}&\prod_{i=1}^g
 \frac{\gamma_i+g-i+1}{g-i+1}\\
 &\times\prod_{1\le i<j\le g}
 \frac{\gamma_i-\gamma_j+j-i}{j-i}
 \frac{\gamma_i+\gamma_j+2g-i-j+2}{2g-i-j+2}.
\end{split}
\label{eq:weyl-dimension}
\end{equation}

\begin{lemma}[Uniform type-\(C\) heat trace]
\label{lem:uniform-heat-trace}
There is an absolute constant \(C\) such that, for \(0<s\le1\),
\begin{equation}
 \boxed{
 K_{2s}(e)=\sum_\gamma d_\gamma^2e^{-2s\kappa_\gamma}
 \le \exp\!\bigl(Cg^2\log(g+2)\bigr)
       s^{-(2g^2+g+1)/2}.}
\label{eq:uniform-heat-trace}
\end{equation}
The identity is for the heat-kernel density relative to Haar probability.
\end{lemma}

\begin{proof}
Put \(R=\gamma_1\).  Equation \eqref{eq:casimir} gives
\(\kappa_\gamma\ge R^2\).  Formula \eqref{eq:weyl-dimension} has exactly
\(g^2\) factors; every denominator is at least one and every numerator is at
most \(2(R+g+1)\).  Hence
\begin{equation}
 d_\gamma\le[2(R+g+1)]^{g^2}.
\label{eq:dimension-crude}
\end{equation}
There are at most \((R+1)^g\) partitions with largest part \(R\), and
\(R+g+1\le(g+1)(R+1)\).  With \(N=2g^2+g\), this gives
\begin{equation}
 K_{2s}(e)
 \le[2(g+1)]^{2g^2}
     \sum_{R=0}^\infty(R+1)^Ne^{-2sR^2}.
\label{eq:heat-shells}
\end{equation}
To bound the last sum, first use \((R+1)^N\le2^NR^N\) for \(R\ge1\).
For a nonnegative unimodal function, its sum over positive integers is at
most its integral plus twice its maximum.  In the present case
\[
 \int_0^\infty x^Ne^{-2sx^2}\,dx
 =\frac12(2s)^{-(N+1)/2}
   \Gamma\!\left(\frac{N+1}{2}\right),
 \qquad
 \max_{x\ge0}x^Ne^{-2sx^2}
 =\left(\frac{N}{4es}\right)^{N/2}.
\]
Stirling's upper bound for the displayed Gamma factor, together with
\(s\le1\), therefore yields
\[
 \sum_{R=0}^\infty(R+1)^Ne^{-2sR^2}
 \le \exp\!\bigl(CN\log(N+2)\bigr)s^{-(N+1)/2}
\]
for \(0<s\le1\).  Substitution into \eqref{eq:heat-shells}, with
\(N=2g^2+g\), proves \eqref{eq:uniform-heat-trace}.
\end{proof}

\subsection{Peter--Weyl smoothing}

The following statement isolates exactly how a character estimate is
converted into a bound for a nonspectral test function.

\begin{lemma}[Brownian displacement]
\label{lem:brownian-displacement}
Let \(M\) be a compact \(D\)-dimensional Riemannian manifold with
nonnegative Ricci curvature, let \(P_s=e^{s\Delta}\) be its heat semigroup,
and let \(K_s(x,dy)\) be the corresponding transition law.  Then
\begin{equation}
 \int_M d(x,y)^2K_s(x,dy)\le2Ds.
\label{eq:brownian-distance}
\end{equation}
\end{lemma}

\begin{proof}
Put \(r(y)=d(x,y)\).  Laplacian comparison gives
\(\Delta r\le(D-1)/r\) off \(x\) and its cut locus, and the same inequality
holds in the barrier, hence distributional, sense.  The chain rule away from
that exceptional set therefore gives
\[
 \Delta r^2=2r\Delta r+2|\nabla r|^2\le2D
\]
distributionally on all of \(M\).  Apply the positive heat semigroup to this
inequality and integrate the identity
\(\partial_uP_ur^2=P_u\Delta r^2\) from \(0\) to \(s\).  Smooth bounded
approximations of \(r^2\) justify the calculation across the cut locus; on a
compact manifold no truncation at infinity is needed.  Since \(r(x)=0\), the
result is \eqref{eq:brownian-distance}.
\end{proof}

\begin{proposition}[Heat bridge]
\label{prop:heat-bridge}
Let \(\mu\) be a conjugation-invariant probability measure on \(G_g\), and
suppose that for every nontrivial irreducible representation \(\rho\),
\begin{equation}
 \left|\int_{G_g}\chi_\rho\,d\mu\right|
 \le\varepsilon d_\rho.
\label{eq:abstract-character-bound}
\end{equation}
If \(h\) is central, \(L\)-Lipschitz, and \(\|h\|_\infty\le1\), then for
\(0<s\le1\),
\begin{equation}
 \boxed{
 \left|\int h\,d\mu-\int h\,dm_g^{\mathrm{Haar}}\right|
 \le2L\sqrt{2D_gs}+\varepsilon K_{2s}(e)^{1/2}.}
\label{eq:abstract-heat-bridge}
\end{equation}
\end{proposition}

\begin{proof}
Under Haar probability, irreducible characters form an orthonormal basis for
the central square-integrable functions.  The Fourier and heat expansions are
\begin{equation}
 h=\sum_\rho\widehat h(\rho)\chi_\rho,
 \qquad
 P_sh=\sum_\rho e^{-s\kappa_\rho}
                  \widehat h(\rho)\chi_\rho,
 \qquad
 \sum_\rho|\widehat h(\rho)|^2=\|h\|_2^2.
\label{eq:PW-expansion}
\end{equation}
The trivial character cancels between \(\mu\) and
\(m_g^{\mathrm{Haar}}\).  Applying
\eqref{eq:abstract-character-bound}, Cauchy--Schwarz, and
\(\|h\|_2\le\|h\|_\infty\) therefore gives
\begin{equation}
 \left|\int P_sh\,d(\mu-m_g^{\mathrm{Haar}})\right|
 \le\varepsilon
 \left(\sum_\rho d_\rho^2e^{-2s\kappa_\rho}\right)^{1/2}
 =\varepsilon K_{2s}(e)^{1/2}.
\label{eq:PW-CS}
\end{equation}
The last equality is exact: relative to Haar probability, an irreducible of
dimension \(d_\rho\) contributes its \(d_\rho^2\) matrix coefficients to the
full heat kernel.  In particular, the quantity in \eqref{eq:PW-CS} is
\(K_{2s}(e)\), not \(K_s(e)\), and no Riemannian-volume factor occurs.

It remains to control smoothing.  A compact group with a bi-invariant metric
has nonnegative Ricci curvature: on its Lie algebra,
\(\operatorname{Ric}(X,X)=\frac14\sum_j\|[X,X_j]\|^2\ge0\).
Thus \cref{lem:brownian-displacement}, translation invariance, and
Cauchy--Schwarz imply
\[
 \|P_sh-h\|_\infty\le L\sqrt{2D_gs}.
\]
Use this once under \(\mu\) and once under \(m_g^{\mathrm{Haar}}\), and insert
\eqref{eq:PW-CS}.  This proves \eqref{eq:abstract-heat-bridge}.
\end{proof}
